% DL9MJ, Matthias Jung, 2026
\begin{tikzpicture}

    % Parameter
    \pgfmathsetmacro{\f}{1.4}       % Brennweite des Mutter-Paraboloids
    \pgfmathsetmacro{\ap}{2.4}      % halbe Apertur
    \pgfmathsetmacro{\ymin}{0.75}   % unterer Rand des Offset-Ausschnitts
    \pgfmathsetmacro{\ymax}{2.55}   % oberer Rand des Offset-Ausschnitts
    \pgfmathsetmacro{\xend}{5.2}    % Ende der reflektierten Strahlen

    % Brennpunkt des Mutter-Paraboloids
    \coordinate (F) at (\f,0);

    % ------------------------------------------------------------
    % Mutter-Paraboloid nur angedeutet
    % ------------------------------------------------------------
    \draw[gray!45, dashed, thick]
        plot[
            variable=\t,
            domain=-2.6:2.6,
            samples=120,
            smooth
        ]
        ({\t*\t/(4*\f)},{\t});

    % Achse des Mutter-Paraboloids
    \draw[gray!45, dashed, thick] (-0.15,0) -- (\xend,0);

    % ------------------------------------------------------------
    % Offsetspiegel: nur ein seitlicher Ausschnitt
    % ------------------------------------------------------------
    \draw[ultra thick]
        plot[
            variable=\t,
            domain=\ymin:\ymax,
            samples=100,
            smooth
        ]
        ({\t*\t/(4*\f)},{\t});

    % Unterer Punkt des Offsetspiegels
    \coordinate (Mlow) at ({\ymin*\ymin/(4*\f)}, {\ymin});

    % ------------------------------------------------------------
    % Feed / Erregerantenne im Brennpunkt
    % ------------------------------------------------------------

    \draw({(-\ap)*(-\ap)/(4*\f)},{-\ap}) coordinate(end);

    % Feed / Antenne im Brennpunkt
    \node[bareantenna, scale=1.0, anchor=west] (ant) at (F) {};

    \draw(end) to [open] ++(2,0) node[trxshape, scale=1, anchor=west](trx){};

    % Anschlussleitung und TRX
    \draw (ant)
        -- ++(0,-1.2) coordinate(h)
        -| (trx.north);

    % ------------------------------------------------------------
    % Strahlen vom Feed zum Offsetspiegel und danach parallel
    % ------------------------------------------------------------
    \foreach \yy in {0.85,1.20,1.55,1.90,2.25,2.50} {
        \pgfmathsetmacro{\xx}{\yy*\yy/(4*\f)}

        % vom Brennpunkt zum Spiegel
        \draw[DARCred, thick]
            (F) -- (\xx,\yy);

        % reflektierter Strahl parallel zur Achse
        \draw[DARCred, thick]
            (\xx,\yy) -- (\xend,\yy);
    }

\end{tikzpicture}